\subsection{关于两测度乘积可测的函数}

命题 2. --- 设 f 是定义在 T×T′ 上、取值于一个拓扑空间 G 中的 ν-可测函数，并设 M 是由使映射 t′ ↦ f(t, t′) 不是 μ′-可测的 t ∈ T 组成之集。

\begin{enumerate}
\def\labelenumi{\alph{enumi})}
\item
  若 \(f\) 在一个 \(\nu\) -适度子集（\(\mathbf{T} \times \mathbf{T}'\) 的子集）的余集上为常数，则 \(\mathbf{M}\) 是 \(\mu\) -可忽略的。
\item
  若 \(\mu'\) 是适度的，则 M 是 \(\mu\) -可忽略的。
\end{enumerate}

断言 a) 由 §3, No.~2 的命题 4b) 以及 No.~1 的诸注记得出。为处理 b)，注意 \(\mu'\) 是有界测度序列 \(\mu_n'\) 之和（§2, No.~3, 命题 4）；\(f\) 关于 \(\mu \otimes \mu_n' \leqslant \nu\) 可测，且集 M 是诸集 \(M_{n}\) 之并，这些集与诸测度 \(\mu_{n}^{\prime}\) 相伴（§2, No.~2, 命题 2）。于是化归到 \(\mu^{\prime}\) 有界的情形，而这由 §3, No.~2 的命题 4c) 得出。

这一陈述立即推广到复测度（Ch. III, §4, No.~2, 命题 3）。

推论. --- 设 A 是一个 \(\nu\) -可测子集（\(T \times T'\) 的子集），并设 M 是由这样的 \(t \in T\) 组成之集：A 在 t 处的截口 \(\mathrm{A}(t)\) 不是 \(\mu'\) -可测。

\begin{enumerate}
\def\labelenumi{\alph{enumi})}
\item
  若 A 是 \(\nu\) -适度的，则 M 是 \(\mu\) -可忽略的。
\item
  若 A 在 \(T'\) 上的投影是 \(\mu'\) -适度的，则 M 是局部 \(\mu\) -可忽略的。
\end{enumerate}

断言 a) 由命题 2 立即得出。为建立 b)，用 B 表示一个集，它是 T′ 中 \(\mu'\) -可积开集序列之并，且包含 A 在 T′ 上的投影，并用 \(\mu_1'\) 表示适度测度 \(\varphi_{\mathrm{B}} \cdot \mu'\) ；由于 A 关于 \(\mu \otimes \mu_1' \leqslant \mu \otimes \mu'\) 可测，命题 2 蕴涵 A(t) 是 \(\mu_1'\) -可测的（除去组成一个局部 \(\mu\) -可忽略集的 t 外）。但因 A(t) ⊂ B，说 A(t) 是 \(\mu_1'\) -可测的就等价于说 A(t) 是 \(\mu'\) -可测的（§5, No.~3, 命题 4 的推论）。

命题 3. --- 设 f 是 T 到一个拓扑空间 F 中的映射。若 f 是 \(\mu\) -可测的，则映射 \((t, t') \mapsto f(t)\) 是 \(\nu\) -可测的。反之，若 \(\mu' \neq 0\) 且此映射是 \(\nu\) -可测的，则函数 f 是 \(\mu\) -可测的。

第一个断言由 §6, No.~6 命题 10 的推论 1 得出。设 \(\mu' \neq 0\) ，用 \(\mu_1'\) 表示一个被 \(\mu'\) 优超的非零具紧支集测度，用 \(\nu_1\) 表示测度 \(\mu \otimes \mu_1'\) ，并置 \(a = \| \mu_1'\|\) 。于是投影 \(\mathrm{pr}_1\) （\(\mathrm{T} \times \mathrm{T}'\) 到 \(\mathrm{T}\) 上的投影）是 \(\nu_1\) -逆紧的，且像测度 \(\mathrm{pr}_1(\nu_1)\) 等于 \(a\mu\) （§6, No.~6, 命题 10）。若 \((t, t') \mapsto f(t)\) 是 \(\nu\) -可测的，则它也是 \(\nu_1\) -可测的，因此 \(f\) 关于测度 \(a\mu\) 可测（§6, No.~2, 命题 3），再由 \(a \neq 0\) 即得结论。

前述陈述立即推广到复测度（Ch. III, §4, No.~2, 命题 3），下列推论亦然。

推论 1. --- 设 F、F′ 和 G 是三个拓扑空间，u 是 F × F′ 到 G 中的一个连续映射。设 f（相应地，f′）是定义在 T（相应地，T′）上、取值于 F（相应地，F′）中且关于 μ（相应地，μ′）可测的函数。则函数 \((t, t') \mapsto u(f(t), f'(t'))\) 关于 μ ⊗ μ′ 可测。

由于映射 \((t,t^{\prime})\mapsto f(t)\) 、\((t,t^{\prime})\mapsto f^{\prime}(t^{\prime})\) 依命题 3 是 \(\nu\) -可测的，这由 Ch. IV, §5, No.~3 的定理 1 得出。

推论 2. --- 若 A ⊂ T 与 A′ ⊂ T′（分别关于 μ 与 μ′）可测，则 A × A′ 关于 μ ⊗ μ′ 可测。

这由推论 1 立即得出。

推论 3. --- 考虑定义在 T 上的正数值（相应地，复值）函数 f 与 \(f'\) （定义在 \(T'\) 上）。若这两个函数分别关于 \(\mu\) 与 \(\mu'\) 可测，则函数

\[
f \otimes f ^ {\prime}: (t, t ^ {\prime}) \mapsto f (t) f ^ {\prime} (t ^ {\prime})
\]

关于 \(\mu \otimes \mu'\) 可测。

复函数或有限实函数的情形是推论 1 的直接推论。为处理正数值函数的情形，对每个整数 \(n \geqslant 0\) 置 \(f_n = \inf(f, n)\) 、\(f_n' = \inf(f', n)\) ，则有（按惯例 \(0 \cdot (+\infty) = 0\) ）\(f \otimes f' = \sup_n(f_n \otimes f_n')\) ，由此即得结论。

命题 4. --- 设 A 是 T 的一个子集。若 A 是局部 \(\mu\) -可忽略的，则 \(A \times T'\) 是局部 \(\nu\) -可忽略的。反之，若 \(A \times T'\) 是局部 \(\nu\) -可忽略的且 \(\mu' \neq 0\) ，则 A 是局部 \(\mu\) -可忽略的。

第一个断言由 §6, No.~6 命题 10 的推论 1 得出。为建立第二个断言，仍用命题 3 证明中的记号；\(A \times T' = \overline{\text{pr}}_{1}(A)\) 对测度 \(\nu_{1}\) 是局部可忽略的，因此 A 对 \(a\mu\) 是局部可忽略的（§6, No.~2, 命题 2 的推论），再由 \(a \neq 0\) 即得结论。

前述陈述立即推广到两个复测度的乘积（Ch. III, §4, No.~2, 命题 3），下列推论亦然。

推论. --- 若测度 \(\mu\) （相应地，\(\mu'\) ）集中于 M（相应地，M′），则 \(\mu \otimes \mu'\) 集中于 M × M′。

因为 \((\mathrm{T} \times \mathrm{T}') - (\mathrm{M} \times \mathrm{M}')\) 是 \((\mathrm{T} - \mathrm{M}) \times \mathrm{T}'\) 与 \(\mathrm{T} \times (\mathrm{T}' - \mathrm{M}')\) 之并，而依命题 4 这两个集对 \(\mu \times \mu'\) 都是局部可忽略的。

